% ============================================================
%  Chapter 5 — Exact equations and integrating factors
% ============================================================
\chapter{Exact Equations and Integrating Factors}\label{ch:ch05}

\begin{diagnostic}
  \item For $F(x,y) = x^2y + \sin(xy)$, compute $\pder{F}{x}$, $\pder{F}{y}$,
        $\pder{}{y}\pder{F}{x}$ and $\pder{}{x}\pder{F}{y}$. What do you notice?
  \item If $y = y(x)$ is differentiable, compute $\der{}{x}F(x,y(x))$ in terms of
        the partial derivatives of $F$.
  \item Evaluate $\displaystyle\der{}{y}\int_0^1 \sin(ty)\,\mathrm{d}t$ in two ways:
        by computing the integral first, and by differentiating inside it.
\end{diagnostic}

% ------------------------------------------------------------
\section{Motivation}\label{sec:ch05-motivation}

Consider
\begin{equation}\label{eq:ch05-motivation}
  \der{y}{x} = -\frac{2xy + 1}{x^2 + 3y^2}.
\end{equation}
It is not separable and not linear. But look at $F(x,y) = x^2y + x + y^3$. Along
any differentiable curve $y = y(x)$,
\[
  \der{}{x}F(x,y(x)) = \pder{F}{x} + \pder{F}{y}\der{y}{x}
  = (2xy + 1) + (x^2 + 3y^2)\der{y}{x},
\]
which is zero exactly when $y$ solves \eqref{eq:ch05-motivation}. So the
solution curves are the level curves $x^2y + x + y^3 = c$. The equation was
secretly the statement ``$F$ is constant''. When does such an $F$ exist, how do
we find it, and what if it does not exist but almost does? Those three
questions are this chapter. The answer to the first is a statement about
mixed partial derivatives, and, surprisingly, about the \emph{shape of the
domain}.

% ------------------------------------------------------------
\section{Exact equations}\label{sec:ch05-exact}

\notation{Leibniz and partial derivatives throughout: the structure of the
differential $\mathrm{d}F = \pder{F}{x}\mathrm{d}x + \pder{F}{y}\mathrm{d}y$ is the
whole point.}

We write first-order equations symmetrically as
\begin{equation}\label{eq:ch05-form}
  M(x,y)\,\mathrm{d}x + N(x,y)\,\mathrm{d}y = 0,
\end{equation}
which means: a curve $y = y(x)$ is a solution if
$M(x,y) + N(x,y)\der{y}{x} = 0$, and a curve $x = x(y)$ is a solution if
$M\der{x}{y} + N = 0$. This form treats $x$ and $y$ equally, which is natural
for level curves (a circle is not a graph over either axis).

\begin{definition}[Exact]\label{def:ch05-exact}
Let $M, N$ be continuous on an open set $D\subseteq\R^2$. Equation
\eqref{eq:ch05-form} is \emph{exact on $D$} if there is a function $F$ with
continuous first partial derivatives on $D$ (a \emph{potential}) such that
\[
  \pder{F}{x} = M,\qquad \pder{F}{y} = N\qquad\text{on } D .
\]
\end{definition}

\begin{theorem}[Solutions of an exact equation]\label{thm:ch05-level}
Suppose \eqref{eq:ch05-form} is exact on $D$ with potential $F$. A
differentiable $\varphi$ on an interval $I$, with $(x,\varphi(x))\in D$, solves
$M + N\der{y}{x} = 0$ iff $F(x,\varphi(x))$ is constant on $I$. Conversely, if
$F(x_0,y_0) = c$ and $N(x_0,y_0)\neq 0$, the level set $F = c$ is, near
$(x_0,y_0)$, the graph of a solution.
\end{theorem}
\begin{proof}
By the chain rule, $\der{}{x}F(x,\varphi(x)) = M(x,\varphi) + N(x,\varphi)\varphi'(x)$.
This vanishes on $I$ iff $\varphi$ is a solution, and a function with zero
derivative on an interval is constant (\cref{thm:ch01-antiderivative}). The
converse is \cref{prop:ch01-implicit} with $G = F - c$ (note
$\pder{G}{y} = N\neq 0$ and $\pder{G}{x} + \pder{G}{y}\cdot(-M/N) = 0$).
\end{proof}

\subsection{The test for exactness}

We need two theorems of multivariable calculus.

\begin{theorem}[Symmetry of mixed partials]\label{thm:ch05-schwarz}
If $F$ has continuous partial derivatives $\pder{F}{x}$, $\pder{F}{y}$,
$\pder{}{y}\pder{F}{x}$ and $\pder{}{x}\pder{F}{y}$ on an open set, then
$\pder{}{y}\pder{F}{x} = \pder{}{x}\pder{F}{y}$ there.
\end{theorem}

\begin{theorem}[Differentiation under the integral sign]\label{thm:ch05-leibniz}
If $f$ and $\pder{f}{x}$ are continuous on $[a,b]\times[c,d]$, then
$\displaystyle\der{}{x}\int_c^d f(x,t)\,\mathrm{d}t = \int_c^d\pder{f}{x}(x,t)\,\mathrm{d}t$
for $a<x<b$.
\end{theorem}
\noindent Both proofs are deferred to \textcite[Ch.~9]{Rudin}.
% VERIFY: Rudin 3rd ed. Thm 9.41 (mixed partials) and Thm 9.42 (differentiation under the integral sign).

\begin{theorem}[Exactness test]\label{thm:ch05-test}
Let $M, N$ have continuous first partial derivatives on an open rectangle
$R = (a,b)\times(c,d)$. Then \eqref{eq:ch05-form} is exact on $R$ iff
\begin{equation}\label{eq:ch05-test}
  \pder{M}{y} = \pder{N}{x}\qquad\text{on } R .
\end{equation}
In that case, for any $(x_0,y_0)\in R$, a potential is
\begin{equation}\label{eq:ch05-potential}
  F(x,y) = \int_{x_0}^{x}M(s,y_0)\,\mathrm{d}s + \int_{y_0}^{y}N(x,t)\,\mathrm{d}t .
\end{equation}
\end{theorem}

\begin{proof}
(\emph{Necessity.}) If $\pder{F}{x} = M$ and $\pder{F}{y} = N$, then the mixed
partials of $F$ are $\pder{M}{y}$ and $\pder{N}{x}$, which are continuous, so
they are equal by \cref{thm:ch05-schwarz}.

(\emph{Sufficiency.}) Define $F$ by \eqref{eq:ch05-potential}; it is defined
because every horizontal and vertical segment used lies in the rectangle.
\begin{align*}
  \pder{F}{y}(x,y) &= 0 + N(x,y)
     \why{FTC in $y$; the first term does not depend on $y$}\\
  \pder{F}{x}(x,y) &= M(x,y_0) + \int_{y_0}^{y}\pder{N}{x}(x,t)\,\mathrm{d}t
     \why{FTC; \cref{thm:ch05-leibniz}}\\
  &= M(x,y_0) + \int_{y_0}^{y}\pder{M}{t}(x,t)\,\mathrm{d}t
     \why{hypothesis \eqref{eq:ch05-test}}\\
  &= M(x,y_0) + \bigl(M(x,y) - M(x,y_0)\bigr) = M(x,y)
     \why{FTC in $t$}.
\end{align*}
Both partials are continuous, so $F$ is a potential.
\end{proof}

The proof used the rectangle in one place: to integrate along the segments from
$(x_0,y_0)$ to $(x,y_0)$ to $(x,y)$, they must stay in the domain. On domains
with holes the test can pass and exactness still fail.

\begin{example}[Needs care: the test passes, exactness fails]\label{ex:ch05-angle}
Let $D = \R^2\setminus\{(0,0)\}$ and
\[
  M = \frac{-y}{x^2+y^2},\qquad N = \frac{x}{x^2+y^2} .
\]
\emph{Test.} By the quotient rule,
$\pder{M}{y} = \frac{-(x^2+y^2) + 2y^2}{(x^2+y^2)^2} = \frac{y^2 - x^2}{(x^2+y^2)^2}$
and $\pder{N}{x} = \frac{(x^2+y^2) - 2x^2}{(x^2+y^2)^2} = \frac{y^2-x^2}{(x^2+y^2)^2}$:
equal on $D$.

\emph{Not exact on $D$.} Suppose $F$ were a potential on $D$. Along the unit
circle $\gamma(\theta) = (\cos\theta,\sin\theta)$, by the chain rule
\begin{align*}
  \der{}{\theta}F(\gamma(\theta))
  &= M\der{x}{\theta} + N\der{y}{\theta}
   = (-\sin\theta)(-\sin\theta) + (\cos\theta)(\cos\theta) = 1,
\end{align*}
so $F(\gamma(2\pi)) - F(\gamma(0)) = 2\pi$. But $\gamma(2\pi) = \gamma(0)$, a
contradiction.

\emph{Exact on smaller domains.} On the half-plane $x>0$ (a region that
contains every segment we need), $F = \arctan(y/x)$ is a potential: it is the
polar angle $\theta$. The angle exists locally but cannot be defined
continuously all the way around the origin. The solution curves of
$-y\,\mathrm{d}x + x\,\mathrm{d}y = 0$ are, as $F = $ const shows, rays from
the origin.
\end{example}

\subsection{Solving an exact equation by hand}

Rather than memorizing \eqref{eq:ch05-potential}, integrate one equation and
correct with the other:
\begin{enumerate}[label=(\arabic*)]
  \item Check $\pder{M}{y} = \pder{N}{x}$.
  \item Integrate $\pder{F}{x} = M$ in $x$, with an unknown function $g(y)$ as
        the ``constant'': $F = \int M\,\mathrm{d}x + g(y)$.
  \item Differentiate in $y$ and set equal to $N$; this determines $g'(y)$.
        (If $g'$ came out depending on $x$, step (1) was wrong.)
  \item The solutions are $F(x,y) = c$; fit $c$ to initial data.
\end{enumerate}

\begin{example}[The motivating equation]\label{ex:ch05-motivation}
Solve $(2xy + 1)\,\mathrm{d}x + (x^2 + 3y^2)\,\mathrm{d}y = 0$, $y(0) = 1$.

\emph{Solution.}
\begin{align*}
  \pder{M}{y} &= 2x = \pder{N}{x} \why{exact on $\R^2$}\\
  F &= \int(2xy + 1)\,\mathrm{d}x = x^2y + x + g(y) \why{$y$ held fixed}\\
  \pder{F}{y} &= x^2 + g'(y) = x^2 + 3y^2 \;\Rightarrow\; g' = 3y^2,\ g = y^3
     \why{match $N$}.
\end{align*}
So $x^2y + x + y^3 = c$, and $(0,1)$ gives $c = 1$:
$x^2y + x + y^3 = 1$. Near $(0,1)$, $N = x^2 + 3y^2 = 3\neq0$, so by
\cref{thm:ch05-level} this defines $y$ as a solution of the IVP near $x = 0$.
\end{example}

\begin{example}[A trigonometric potential]\label{ex:ch05-trig}
Solve $(3x^2 + y\cos x)\,\mathrm{d}x + (\sin x - 4y^3)\,\mathrm{d}y = 0$.

\emph{Solution.}
\begin{align*}
  \pder{M}{y} &= \cos x = \pder{N}{x} \why{exact}\\
  F &= x^3 + y\sin x + g(y) \why{integrate $M$ in $x$}\\
  \pder{F}{y} &= \sin x + g'(y) = \sin x - 4y^3 \;\Rightarrow\; g = -y^4 .
\end{align*}
Solutions: $x^3 + y\sin x - y^4 = c$.
\end{example}

% ------------------------------------------------------------
\section{Integrating factors}\label{sec:ch05-integrating}

\begin{definition}\label{def:ch05-if}
A nonvanishing function $\mu$ with continuous first partials is an
\emph{integrating factor} for \eqref{eq:ch05-form} on $D$ if
$\mu M\,\mathrm{d}x + \mu N\,\mathrm{d}y = 0$ is exact on $D$.
\end{definition}

Multiplying by a nonvanishing $\mu$ does not change which curves are
solutions (both equations say $M + N\der{y}{x} = 0$). On a rectangle, by
\cref{thm:ch05-test}, $\mu$ is an integrating factor iff
\begin{equation}\label{eq:ch05-if-pde}
  \pder{(\mu M)}{y} = \pder{(\mu N)}{x}
  \quad\Longleftrightarrow\quad
  M\pder{\mu}{y} - N\pder{\mu}{x} = \mu\Bigl(\pder{N}{x} - \pder{M}{y}\Bigr).
\end{equation}
This is a \emph{partial} differential equation for $\mu$, generally harder than
the original problem. We look for $\mu$ depending on one variable only.

\begin{theorem}[One-variable integrating factors]\label{thm:ch05-if-special}
Let $M,N$ have continuous first partials on a rectangle $R$.
\begin{enumerate}[label=(\alph*)]
  \item If $N\neq 0$ and $\dfrac{\pder{M}{y} - \pder{N}{x}}{N} = \rho(x)$ depends
        on $x$ only, then $\mu(x) = e^{\int\rho(x)\,\mathrm{d}x}$ is an integrating
        factor.
  \item If $M\neq0$ and $\dfrac{\pder{N}{x} - \pder{M}{y}}{M} = \sigma(y)$ depends
        on $y$ only, then $\mu(y) = e^{\int\sigma(y)\,\mathrm{d}y}$ is an integrating
        factor.
\end{enumerate}
\end{theorem}
\begin{proof}
(a) For $\mu = \mu(x)$, $\pder{\mu}{y} = 0$, so \eqref{eq:ch05-if-pde} reads
$-N\der{\mu}{x} = \mu\bigl(\pder{N}{x} - \pder{M}{y}\bigr)$, i.e.
$\der{\mu}{x} = \rho(x)\mu$. The given $\mu$ solves this (it is positive, and
$\der{\mu}{x} = \rho\mu$ by the chain rule). (b) is the same computation with
the roles of $x$ and $y$ exchanged.
\end{proof}

\begin{example}[Linear equations, revisited]\label{ex:ch05-linear}
Write $\der{y}{x} + p(x)y = q(x)$ as $\bigl(p(x)y - q(x)\bigr)\mathrm{d}x + \mathrm{d}y = 0$.
Then $\frac{\pder{M}{y} - \pder{N}{x}}{N} = \frac{p - 0}{1} = p(x)$, so
$\mu = e^{\int p}$: \cref{thm:ch05-if-special}(a) \emph{is} the integrating
factor of Chapter~4.
\end{example}

\begin{example}[An integrating factor in $x$]\label{ex:ch05-mu-x}
Solve $(x^2 + y^2 + x)\,\mathrm{d}x + xy\,\mathrm{d}y = 0$.

\emph{Solution.}
\begin{align*}
  \pder{M}{y} &= 2y,\quad \pder{N}{x} = y \why{not exact}\\
  \frac{\pder{M}{y} - \pder{N}{x}}{N} &= \frac{y}{xy} = \frac1x
     \why{depends on $x$ only, for $xy\neq 0$}\\
  \mu &= e^{\ln|x|} = |x|;\ \text{take } \mu = x \why{a constant multiple is fine}\\
  (x^3 + xy^2 + x^2)\,\mathrm{d}x &+ x^2y\,\mathrm{d}y = 0
     \why{check: both mixed partials are $2xy$}\\
  F &= \frac{x^4}{4} + \frac{x^2y^2}{2} + \frac{x^3}{3}
     \why{integrate in $x$; then $g'(y) = 0$}.
\end{align*}
Solutions: $3x^4 + 6x^2y^2 + 4x^3 = c$. \emph{Care:} $\mu = x$ vanishes on the
line $x = 0$. That line is itself a solution curve of the original form
($\mathrm{d}x = 0$ along it, and $M\cdot 0 + N\,\mathrm{d}y = 0$ since $N = 0$
there), and it is the level set $c = 0$ of $F$ only by accident. Whenever $\mu$
vanishes or blows up, check those curves separately.
\Cref{fig:ch05-levels} shows the level curves.
\end{example}

\begin{figure}[ht]
  \centering
  % Level curves x^4/4 + x^2 y^2/2 + x^3/3 = c, i.e.
% y = +- sqrt(2(c - x^4/4 - x^3/3)) / |x|.
\begin{tikzpicture}
\begin{axis}[deplot, xmin=-2, xmax=1.6, ymin=-3, ymax=3,
  xlabel={$x$}, ylabel={$y$}, unbounded coords=jump,
  restrict y to domain=-3.2:3.2]
  \pgfplotsinvokeforeach{-0.06,-0.02,0.05,0.3,1}{
    \addplot[blue, samples=400, domain=-2:1.6]
      {sqrt(2*(#1 - x^4/4 - x^3/3))/abs(x)};
    \addplot[blue, samples=400, domain=-2:1.6]
      {-sqrt(2*(#1 - x^4/4 - x^3/3))/abs(x)};
  }
  \addplot[red!70!black, thick] coordinates {(0,-3) (0,3)};
\end{axis}
\end{tikzpicture}

  \caption{Solution curves $\tfrac{x^4}{4} + \tfrac{x^2y^2}{2} + \tfrac{x^3}{3} = c$
  of \cref{ex:ch05-mu-x} for $c\in\{-0.06,-0.02,0.05,0.3,1\}$. The
  $y$-axis is also a solution curve.}
  \label{fig:ch05-levels}
\end{figure}

\begin{example}[An integrating factor in $y$]\label{ex:ch05-mu-y}
Solve $y\,\mathrm{d}x + (2x - ye^y)\,\mathrm{d}y = 0$.

\emph{Solution.}
\begin{align*}
  \frac{\pder{N}{x} - \pder{M}{y}}{M} &= \frac{2 - 1}{y} = \frac1y
     \;\Rightarrow\;\mu = y \why{\cref{thm:ch05-if-special}(b)}\\
  y^2\,\mathrm{d}x &+ (2xy - y^2e^y)\,\mathrm{d}y = 0 \why{mixed partials $2y$}\\
  F &= xy^2 + g(y),\quad 2xy + g'(y) = 2xy - y^2e^y\\
  g &= -\int y^2e^y\,\mathrm{d}y = -(y^2 - 2y + 2)e^y \why{parts twice}.
\end{align*}
Solutions: $xy^2 - (y^2 - 2y + 2)e^y = c$, together with $y\equiv 0$ (where
$\mu$ vanishes; check: $M = 0$ and $\mathrm{d}y = 0$ there).
\end{example}

\begin{example}[Needs care: many integrating factors, different losses]\label{ex:ch05-many}
For $y\,\mathrm{d}x - x\,\mathrm{d}y = 0$, each of
$\mu = \frac1{x^2}$, $\frac1{y^2}$, $\frac{1}{xy}$, $\frac{1}{x^2+y^2}$ is an
integrating factor on a suitable region:
\[
  \frac{y\,\mathrm{d}x - x\,\mathrm{d}y}{x^2} = -\mathrm{d}\Bigl(\frac yx\Bigr),\quad
  \frac{y\,\mathrm{d}x - x\,\mathrm{d}y}{y^2} = \mathrm{d}\Bigl(\frac xy\Bigr),\quad
  \frac{y\,\mathrm{d}x - x\,\mathrm{d}y}{xy} = \mathrm{d}\ln\Bigl|\frac xy\Bigr| .
\]
All describe lines through the origin, but $\mu = 1/x^2$ is undefined on $x = 0$
and so its potential $y/x$ cannot produce the solution $x = 0$; $\mu = 1/y^2$
loses $y = 0$; $\mu = 1/(xy)$ loses both. The region of validity of an
integrating factor is part of the answer. (\Cref{prob:ch05-h3} explains why the
ratio of two integrating factors is always constant along solution curves.)
\end{example}

% ------------------------------------------------------------
\section{Pitfalls}\label{sec:ch05-pitfalls}

\begin{pitfall}[Testing the wrong partials]\label{pit:ch05-wrong-test}
The test is $\pder{M}{y} = \pder{N}{x}$, ``cross'' partials. Checking
$\pder{M}{x} = \pder{N}{y}$ tests nothing. Mnemonic: $M$ goes with
$\mathrm{d}x$, so its \emph{other} variable is $y$.
\end{pitfall}

\begin{pitfall}[A constant instead of a function]\label{pit:ch05-constant}
When integrating $\pder{F}{x} = M$ in $x$, the ``constant'' may depend on $y$:
write $+\,g(y)$, never $+\,C$. Otherwise you will find an $F$ whose
$y$-derivative does not equal $N$.
\end{pitfall}

\begin{pitfall}[Double counting]\label{pit:ch05-double}
Integrating $M$ in $x$ and $N$ in $y$ separately and \emph{adding} the results
counts the mixed terms twice: for $2xy\,\mathrm{d}x + x^2\,\mathrm{d}y$, adding
gives $2x^2y$, but $F = x^2y$. Use the procedure, or take the union of terms
without repetition and then \emph{check} by differentiating.
\end{pitfall}

\begin{pitfall}[Ignoring the domain]\label{pit:ch05-domain}
$\pder{M}{y} = \pder{N}{x}$ guarantees a potential on rectangles (more
generally on star-shaped regions: \cref{prob:ch05-h2}), not on domains with
holes (\cref{ex:ch05-angle}).
\end{pitfall}

\begin{pitfall}[Forgetting where $\mu$ vanishes]\label{pit:ch05-mu-zero}
An integrating factor can add spurious curves (where $\mu = 0$) or hide genuine
ones (where $\mu$ is undefined). Check those curves in the \emph{original}
equation.
\end{pitfall}

% ------------------------------------------------------------
\section{Problems}\label{sec:ch05-problems}

\subsection*{Warm-up}

\begin{problem}{W}\label{prob:ch05-w1}
Which are exact on $\R^2$?
(a) $(2x + y)\,\mathrm{d}x + (x - 2y)\,\mathrm{d}y = 0$;
(b) $y^2\,\mathrm{d}x + (2xy + 1)\,\mathrm{d}y = 0$;
(c) $(x - y)\,\mathrm{d}x + (x + y)\,\mathrm{d}y = 0$.
\end{problem}

\begin{problem}{W}\label{prob:ch05-w2}
Solve $(2x + y)\,\mathrm{d}x + (x - 2y)\,\mathrm{d}y = 0$, $y(1) = 0$.
\end{problem}

\begin{problem}{W}\label{prob:ch05-w3}
Solve $(e^x\sin y + 2x)\,\mathrm{d}x + e^x\cos y\,\mathrm{d}y = 0$.
\end{problem}

\begin{problem}{W}\label{prob:ch05-w4}
Find $k$ so that $(xy^2 + kx^2y)\,\mathrm{d}x + (x^3 + x^2y)\,\mathrm{d}y = 0$ is
exact, and solve it.
\end{problem}

\begin{problem}{W}\label{prob:ch05-w5}
Show that a separable equation $\der{y}{x} = g(x)h(y)$ becomes exact after
division by $h(y)$, and relate the potential to $H$ and $G$ of
\cref{thm:ch03-separation}.
\end{problem}

\subsection*{Core}

\begin{problem}{C}\label{prob:ch05-c1}
Solve $(y\cos x + 2xe^y)\,\mathrm{d}x + (\sin x + x^2e^y - 1)\,\mathrm{d}y = 0$,
$y(0) = 0$.
\end{problem}

\begin{problem}{C}\label{prob:ch05-c2}
Find an integrating factor and solve $(3xy + y^2)\,\mathrm{d}x + (x^2 + xy)\,\mathrm{d}y = 0$.
List any solution curves the integrating factor might add or hide.
\end{problem}

\begin{problem}{C}\label{prob:ch05-c3}
Find an integrating factor and solve $y\,\mathrm{d}x + (y^2 - x)\,\mathrm{d}y = 0$.
\end{problem}

\begin{problem}{C}\label{prob:ch05-c4}
Starting from \cref{ex:ch05-linear}, compute the potential $F$ of
$e^{P(x)}\bigl[(p y - q)\,\mathrm{d}x + \mathrm{d}y\bigr]$ and show that $F = c$
reproduces formula \eqref{eq:ch04-formula}.
\end{problem}

\begin{problem}{C}\label{prob:ch05-c5}
Solve $(y^2 + 2x)\,\mathrm{d}x + 2xy\,\mathrm{d}y = 0$, $y(1) = 1$, explicitly,
and find the maximal interval of the solution $y(x)$.
\end{problem}

\begin{problem}{C}\label{prob:ch05-c6}
For the form of \cref{ex:ch05-angle}, compute $\int_\gamma M\,\mathrm{d}x + N\,\mathrm{d}y$
along the circle of radius $2$ centred at $(3,0)$, and explain the answer using
the potential on $x>0$.
\end{problem}

\begin{problem}{C}\label{prob:ch05-c7}
Find $m, n$ such that $x^my^n$ is an integrating factor of
$(y^2 + 2xy)\,\mathrm{d}x - x^2\,\mathrm{d}y = 0$, solve the equation, and check
whether $y\equiv0$ is a solution the integrating factor hides.
\end{problem}

\begin{problem}{C}\label{prob:ch05-c8}
Show that $\mu = \dfrac{1}{x^2+y^2}$ is an integrating factor of
$(x^2 + y^2 + y)\,\mathrm{d}x - x\,\mathrm{d}y = 0$ on $x>0$, and solve the
equation explicitly for $y$. (Recall \cref{ex:ch05-angle}.)
\end{problem}

\subsection*{Hard}

\begin{problem}{H}\label{prob:ch05-h1}
Show that \eqref{eq:ch05-form} has an integrating factor of the form
$\mu = \phi(xy)$ if $\dfrac{\pder{M}{y} - \pder{N}{x}}{yN - xM}$ is a function of
$z = xy$ alone, and find $\phi$. Apply this to
$y(1 + xy)\,\mathrm{d}x + x(1 - xy)\,\mathrm{d}y = 0$.
\end{problem}

\begin{problem}{H}\label{prob:ch05-h2}
(Poincar\'e lemma in the plane.) Let $D$ be open and \emph{star-shaped} about
the origin ($(x,y)\in D\Rightarrow (tx,ty)\in D$ for $0\le t\le1$), and let
$M, N$ have continuous first partials on $D$ with $\pder{M}{y} = \pder{N}{x}$.
Prove that
\[
  F(x,y) = \int_0^1\bigl[xM(tx,ty) + yN(tx,ty)\bigr]\,\mathrm{d}t
\]
is a potential. Why does this not contradict \cref{ex:ch05-angle}?
\end{problem}

\begin{problem}{H}\label{prob:ch05-h3}
Let $\mu_1,\mu_2$ be integrating factors of \eqref{eq:ch05-form} on a rectangle,
and $F_1$ a potential for $\mu_1M\,\mathrm{d}x + \mu_1N\,\mathrm{d}y$, with
$\nabla F_1\neq 0$. Prove that $g = \mu_2/\mu_1$ is constant along every
solution curve. Check this for $\mu_1 = 1/x^2$, $\mu_2 = 1/y^2$ in
\cref{ex:ch05-many}.
\end{problem}

\begin{problem}{H}\label{prob:ch05-h4}
(Exact second-order equations.) \notation{Lagrange, for the algebra of
coefficients.} Prove that
$P(x)y'' + Q(x)y' + R(x)y = 0$ (with $P\in C^2$, $Q\in C^1$) can be written as
$\der{}{x}\bigl[P y' + (Q - P')y\bigr] = 0$ iff $P'' - Q' + R = 0$. Use this
to find the general solution of $x^2y'' + 3xy' + y = 0$ on $x>0$.
\end{problem}

\begin{problem}{H}\label{prob:ch05-h5}
Show that \eqref{eq:ch05-form} has an integrating factor of the form
$\mu = \phi(x^2 + y^2)$ if
$\dfrac{\pder{M}{y} - \pder{N}{x}}{2(xN - yM)}$ is a function of $z = x^2 + y^2$
alone. Explain how this criterion produces the integrating factor of
\cref{prob:ch05-c8}.
\end{problem}

\subsection*{Putnam/olympiad-style}

\begin{problem}{P}\label{prob:ch05-p1}
Let $M, N$ have continuous first partials on $\R^2$ with
$\pder{M}{y} = \pder{N}{x}$ and $N(x,y)\ge 1$ everywhere. Prove that through
every point of the plane passes exactly one solution curve of
$M\,\mathrm{d}x + N\,\mathrm{d}y = 0$, and that it is the graph of a $C^1$
function defined on all of $\R$.
\end{problem}

\begin{problem}{P}\label{prob:ch05-p2}
Let $P, Q$ be real polynomials in $x,y$ with $\pder{P}{y} = \pder{Q}{x}$. Prove
that $P\,\mathrm{d}x + Q\,\mathrm{d}y$ has a polynomial potential $F$ with
$\deg F\le 1 + \max(\deg P,\deg Q)$, and that the bound can be attained.
\end{problem}

\begin{problem}{P}\label{prob:ch05-p3}
Let $\gamma:[0,1]\to\R^2\setminus\{(0,0)\}$ be a $C^1$ closed curve. Prove that
$\displaystyle\int_\gamma\frac{x\,\mathrm{d}y - y\,\mathrm{d}x}{x^2+y^2}\in 2\pi\mathbb{Z}$.
\end{problem}

% ------------------------------------------------------------
\section{Hints}\label{sec:ch05-hints}

\paragraph{Warm-up and core.}
\cref{prob:ch05-w5}: divide by $h(y)$ and compare with \eqref{eq:ch03-HG}.
\cref{prob:ch05-c2}: try \cref{thm:ch05-if-special}(a).
\cref{prob:ch05-c3}: try \cref{thm:ch05-if-special}(b).
\cref{prob:ch05-c5}: the level set is quadratic in $y$; which sign does the
initial condition pick, and where does the radicand vanish?
\cref{prob:ch05-c6}: the circle lies in $x>0$.
\cref{prob:ch05-c7}: substitute $\mu = x^my^n$ into \eqref{eq:ch05-if-pde} and
match monomials.
\cref{prob:ch05-c8}: write the equation as
$\mathrm{d}x + \frac{y\,\mathrm{d}x - x\,\mathrm{d}y}{x^2+y^2} = 0$.

\hintfor{prob:ch05-h1}
\begin{hintlevels}
  \item Substitute $\mu = \phi(xy)$ into \eqref{eq:ch05-if-pde}; use
        $\pder{\mu}{x} = y\phi'$ and $\pder{\mu}{y} = x\phi'$.
  \item You get $\frac{\phi'(z)}{\phi(z)} = $ (the given ratio).
  \item For the example the ratio is $-2/(xy)$.
\end{hintlevels}

\hintfor{prob:ch05-h2}
\begin{hintlevels}
  \item Differentiate under the integral sign (\cref{thm:ch05-leibniz}).
  \item Use $\pder{M}{y} = \pder{N}{x}$ to recognize the integrand of
        $\pder{F}{x}$ as $\der{}{t}\bigl[tM(tx,ty)\bigr]$.
  \item The punctured plane is not star-shaped about any point.
\end{hintlevels}

\hintfor{prob:ch05-h3}
\begin{hintlevels}
  \item $\mu_2M\,\mathrm{d}x + \mu_2N\,\mathrm{d}y = g\,\mathrm{d}F_1$ is exact.
  \item Exactness of $gF_{1,x}\,\mathrm{d}x + gF_{1,y}\,\mathrm{d}y$ gives
        $g_yF_{1,x} = g_xF_{1,y}$.
  \item The tangent to a solution curve is $(F_{1,y}, -F_{1,x})$; differentiate
        $g$ in that direction.
\end{hintlevels}

\hintfor{prob:ch05-h4}
\begin{hintlevels}
  \item Expand $\der{}{x}[Py' + (Q-P')y]$ and compare coefficients.
  \item For $x^2y'' + 3xy' + y$: $P'' - Q' + R = 2 - 3 + 1 = 0$.
  \item The first integral is a \emph{linear first-order} equation (Chapter~4).
\end{hintlevels}

\hintfor{prob:ch05-h5}
\begin{hintlevels}
  \item With $\mu = \phi(z)$: $\pder{\mu}{x} = 2x\phi'$, $\pder{\mu}{y} = 2y\phi'$.
  \item Substitute into \eqref{eq:ch05-if-pde} and divide.
  \item For \cref{prob:ch05-c8}'s equation, factor $xN - yM$.
\end{hintlevels}

\hintfor{prob:ch05-p1}
\begin{hintlevels}
  \item Build a global potential (the plane is a rectangle).
  \item For fixed $x$, the map $y\mapsto F(x,y)$ is strictly increasing with
        derivative $\ge 1$. What is its range?
  \item Each level set meets every vertical line exactly once; use
        \cref{thm:ch00-ift} for smoothness.
\end{hintlevels}

\hintfor{prob:ch05-p2}
\begin{hintlevels}
  \item Use formula \eqref{eq:ch05-potential} with $(x_0,y_0) = (0,0)$.
  \item Integrating a polynomial in one variable raises degree by at most one.
  \item For attainment, try $P = x^n$, $Q = 0$.
\end{hintlevels}

\hintfor{prob:ch05-p3}
\begin{hintlevels}
  \item Construct a continuous angle along $\gamma$ by solving an ODE:
        $\theta' = \frac{xy' - yx'}{x^2+y^2}$.
  \item With $r = \sqrt{x^2+y^2}$ and $\theta(0)$ chosen so the claim holds at
        $t = 0$, show $(x + iy)/(re^{i\theta})$ has zero derivative; hence
        $x + iy = re^{i\theta}$ all along the curve.
  \item The integral equals $\theta(1) - \theta(0)$, and
        $e^{i\theta(1)} = e^{i\theta(0)}$.
\end{hintlevels}

% ------------------------------------------------------------
\section{Further reading}\label{sec:ch05-reading}

\begin{itemize}
  \item \textcite[\S2.6]{BoyceDiPrima}: exact equations and integrating factors.
  \item \textcite[Lessons 9--10]{TenenbaumPollard}: exact equations, recognizable
        exact differentials, and integrating factors, with many examples of
        ``inspection'' factors like those in \cref{ex:ch05-many}.
  % VERIFY: Tenenbaum & Pollard Lesson 10 title (recognizable exact equations; integrating factors).
  \item \textcite[Ch.~10]{Rudin}: differential forms, closed versus exact forms,
        and the general Poincar\'e lemma behind \cref{ex:ch05-angle} and
        \cref{prob:ch05-h2}.
\end{itemize}
